\Needspace{12\baselineskip}
\section{Smooth nonuniqueness for a degenerate density}
\label{sec:degenerate-hermitian-nonuniqueness}

The first example puts the freedom in the potential inside an open region where the density is zero. Three independent rank-one directions reduce equality of the Monge--Amp\`ere measures to equality of scalar products. We construct two semipositive forms with the same cube, then add an exact correction to obtain a strictly positive equation background.

Let $X=\mathbb P(\mathbb C^2\oplus\mathbb C^2)=\mathbb P^3$, with block homogeneous coordinates $[A:B]$, where $A,B\in\mathbb C^2$, and let $\omega_X$ be its Fubini--Study form.

\begin{theorem}\label{thm:degenerate-hermitian-nonuniqueness}
On $X=\mathbb P^3$ there exist a smooth Hermitian form $\omega$ with $\mathrm{d}\omega\ne0$, a smooth nonnegative nonzero function $f$, and distinct smooth functions $u,v\in\PSH(X,\omega)$ such that
\[
 \sup_Xu=\sup_Xv=0,\qquad
 (\omega+\ddc u)^3=f\omega_X^3=(\omega+\ddc v)^3.
\]
The density vanishes on a nonempty open set. Both potentials have minimal singularities.
\end{theorem}

\subsection{Three rank-one directions}
\label{subsec:numa-rank-directions}

On $U=\{A\ne0,B\ne0\}$ put
\[
 r=\log\frac{|A|^2}{|B|^2},\qquad
 P=\ddc\log|A|^2,\quad Q=\ddc\log|B|^2,\quad
 R=\mathrm{d}r\wedge \dc r.
\]
The logarithms individually use a local homogeneous lift; their Hessians are independent of that lift. The forms $P,Q$ are pullbacks of Fubini--Study forms from the two projective lines. All three forms are semipositive of rank one, and
\begin{equation}\label{eq:numa-rank-frame}
 P^2=Q^2=R^2=0,\qquad P\wedge Q\wedge R>0.
\end{equation}
Indeed, in local coordinates $A=(1,a)$ and $B=\tau(1,b)$, $P$ and $Q$ are positive multiples of $\mathrm{i}\,\mathrm{d}a\wedge \mathrm{d}\bar a$ and $\mathrm{i}\,\mathrm{d}b\wedge \mathrm{d}\bar b$, while
\[
 \partial r=\frac{\bar a\,\mathrm{d}a}{1+|a|^2}
             -\frac{\mathrm{d}\tau}{\tau}
             -\frac{\bar b\,\mathrm{d}b}{1+|b|^2}.
\]
\Needspace{6\baselineskip}
Its $\mathrm{d}\tau$ component is nonzero. The three rank-one directions are therefore independent. In particular, for nonnegative $x,y,z$,
\begin{equation}\label{eq:numa-cube}
 (xP+yQ+zR)^3=6xyz\,P\wedge Q\wedge R,
\end{equation}
and the form is positive exactly when all three coefficients are positive.

\Needspace{5\baselineskip}
For a smooth function $G:\mathbb R\to\mathbb R$ the chain rule is
\begin{equation}\label{eq:numa-chain}
 \ddc G(r)=G'(r)(P-Q)+G''(r)R.
\end{equation}
\Needspace{7\baselineskip}
Let $\sigma(r)=\mathrm{e}^r/(1+\mathrm{e}^r)$. Since
$\log(|A|^2+|B|^2)=\log|B|^2+\log(1+\mathrm{e}^r)$, we have
\begin{equation}\label{eq:numa-fs}
 \omega_X=\sigma(r)P+\sigma(-r)Q+\sigma'(r)R.
\end{equation}
A function of $r$ constant at each end of $\mathbb R$ extends smoothly over $X\setminus U$. A form whose coefficients equal those of \eqref{eq:numa-fs} for $|r|\geq2$ extends smoothly by $\omega_X$ there.

\subsection{The scalar coefficients}
\label{subsec:numa-scalar-coefficients}

The central coefficients will make a different rank-one factor vanish in each solution form. Outside the central interval they will return to the Fubini--Study coefficients, so the resulting forms extend across the omitted projective lines. Primes denote derivatives with respect to $r$.

Choose an even smooth bump $\beta$ supported on $[-1,1]$, strictly positive on $(-1,1)$ and with $\beta'>0$ on $(-1,0)$, normalized by $\int_{-1}^1\beta=2$. A positive multiple of $\exp(-1/(1-r^2))$ on $(-1,1)$, extended by zero, does this. Define
\[
 F(r)=-1+\int_{-1}^r\beta(t)\,\mathrm{d}t\quad(-1<r<1),
 \qquad F=-1\ (r\leq-1),\quad F=1\ (r\geq1),
\]
and put $p=\beta/2$, $q=\beta'/2$. Then $F$ is smooth and odd, $p$ is even, and $q$ is odd, positive on $(-1,0)$ and negative on $(0,1)$. Both $p$ and $q$ vanish to infinite order at the endpoints.

Take smooth $a\geq0$ with $a=1$ on $(-\infty,-2/3]$ and $a=0$ on $[-1/3,\infty)$, and put $b(r)=a(-r)$. Choose an even nonnegative $\chi\in C_c^\infty((-1/3,1/3))$ with $\chi(0)>0$. On $[-1,1]$ set
\[
 z_0=(q^2+\chi^2)^{1/2}.
\]
It is smooth: the expression under the square root is positive in the interior; near each endpoint $z_0=|q|$ with fixed sign of $q$. It is flat at the endpoints and positive in the interior. Also,
\begin{equation}\label{eq:numa-coefficient-cancellation}
 z_0\geq|q|,\qquad a(z_0-q)=0,\qquad b(z_0+q)=0.
\end{equation}
The cancellations follow from the disjoint supports of $a$ and $\chi$, and from the sign of $q$ on the support of $a$; the other equality follows by reflection.

\Needspace{12\baselineskip}
On $(-\infty,-1]$ choose smooth functions $\ell,s$ such that $0\leq\ell\leq1$, $\ell=0$ on $(-\infty,-2]$, $\ell=1$ on $[-3/2,-1]$, while $s>0$ for $r<-1$, $s=1$ on $(-\infty,-3/2]$, and $s$ is flat with value zero at $-1$. Such functions are obtained from ordinary smooth cutoffs and a flat positive endpoint function. Define
\[
 \begin{aligned}
 x&=(1-\ell)\sigma+\ell,&
 y&=(1-\ell)\sigma(-r)+\ell s,&
 z&=(1-\ell)\sigma'+\ell s &&(r\leq-1),\\
 x&=p+a,& y&=p+b,& z&=z_0 &&(-1\leq r\leq1),
 \end{aligned}
\]
and for $r\geq1$ use $x(r)=y(-r)$, $y(r)=x(-r)$, $z(r)=z(-r)$.

The definitions join smoothly. At $-1$ their values are $(1,0,0)$, with all discrepancies flat; near that endpoint $a=1$ and $b=0$. The central definitions are symmetric under reflection, so the same holds at $1$. We have
\[
 x,y,z>0\quad\hbox{for }r\ne\pm1,\qquad
 (x,y,z)(-1)=(1,0,0),\quad (x,y,z)(1)=(0,1,0).
\]
For $|r|\geq2$ their values are $(\sigma(r),\sigma(-r),\sigma'(r))$.

\subsection{Two forms with the same cube}
\label{subsec:numa-equal-cubes}

The cancellation identities now turn the scalar coefficients into two semipositive forms with the same cube. Their average will be strictly positive except at the two ends of the central interval.

The globally smooth form and function
\[
 \omega_0=x(r)P+y(r)Q+z(r)R,\qquad w=F(r)
\]
are understood by the extensions above. The form is semipositive, and positive except on $M_- =\{r=-1\}$ and $M_+=\{r=1\}$. Put
\[
 H=\tfrac12\ddc w=p(P-Q)+qR,\qquad A_0=\omega_0-H,\quad B_0=\omega_0+H.
\]
For $|r|\geq1$ we have $H=0$. On $|r|<1$,
\[
 A_0=aP+(2p+b)Q+(z_0-q)R,\qquad
 B_0=(2p+a)P+bQ+(z_0+q)R.
\]
All coefficients are nonnegative. The cancellations in \eqref{eq:numa-coefficient-cancellation} give $A_0^3=B_0^3=0$ on this band. Outside it both cubes equal $\omega_0^3$. Thus $A_0,B_0$ are globally smooth semipositive forms with equal cubes. That common top-degree form is smooth, vanishes on the closed band $|r|\leq1$, and is positive on its complement, including near $X\setminus U$.

\Needspace{12\baselineskip}
\subsection{Restoring strict positivity of the background}
\label{subsec:numa-positive-background}

The forms $A_0,B_0$ already have the required common measure. We obtain a Hermitian background by adding the same exact correction to their average and subtracting its potential from both solutions. The correction must restore positivity at $M_\pm$ while preserving it in the transition regions.

Take $\kappa\in C_c^\infty((-2,2))$, equal to one on $[-3/2,3/2]$. For a parameter $\varepsilon>0$ to be chosen, set
\[
 \rho(r)=\varepsilon(r^2-1)\kappa(r),\qquad
 \omega=\omega_0+\ddc\rho.
\]
Its coefficients on $U$ are $x+\rho'$, $y-\rho'$ and $z+\rho''$. On $[-3/2,3/2]$ one has $\rho'=2\varepsilon r$ and $\rho''=2\varepsilon$. The last coefficient is positive. For $r\geq0$, $x+2\varepsilon r$ is positive, including at $r=0$ since $x(0)=p(0)>0$. On $[-3/2,0]$, the function $x$ has a positive minimum $m$, so $x+2\varepsilon r\geq m-3\varepsilon>0$ for small $\varepsilon$. Reflection gives the corresponding estimate for $y-2\varepsilon r$. On the compact region $3/2\leq|r|\leq2$, all three original coefficients have a positive minimum and the perturbations are $O(\varepsilon)$. Outside that region and the central interval the form is exactly $\omega_X$. This proves strict positivity globally.

For completeness, nonclosedness can be checked directly. Since $\mathrm{d}R=-\mathrm{d}r\wedge(P-Q)$,
\begin{equation}\label{eq:numa-differential}
 \mathrm{d}\omega=\mathrm{d}\omega_0
   =\mathrm{d}r\wedge\bigl((x'-z)P+(y'+z)Q\bigr).
\end{equation}
At $r=0$, $x'=y'=0$ and $z=\chi(0)>0$. Hence $\mathrm{d}\omega=z\,\mathrm{d}r\wedge(Q-P)\ne0$ there. The nonzero $\mathrm{d}\tau$ component used in \eqref{eq:numa-rank-frame} also shows that $\mathrm{d}r\wedge P$ and $\mathrm{d}r\wedge Q$ are linearly independent.

\begin{proof}[of \cref{thm:degenerate-hermitian-nonuniqueness}]
Set $\widetilde u=-w/2-\rho$ and $\widetilde v=w/2-\rho$. The construction gives
\[
 \omega+\ddc\widetilde u=A_0\geq0,\qquad
 \omega+\ddc\widetilde v=B_0\geq0.
\]
Subtract their respective suprema to obtain $u,v$. Their Hessians are unchanged, and $v-u=w+\text{constant}$ is nonconstant. The function $f=A_0^3/\omega_X^3=B_0^3/\omega_X^3$ is smooth, nonnegative and nonzero. It vanishes on the nonempty open band $|r|<1$. Both smooth potentials have minimal singularities because $V_\omega=0$ for a Hermitian background.
\end{proof}

Every Hermitian background on $\mathbb P^3$ has bounded mass: the Fubini--Study form has invariant mass by closed bounded Stokes, and \cref{prop:bmp-metric} makes finiteness independent of the metric. This applies to all the projective backgrounds constructed below.

The linearized difference also explains the failure of the smooth maximum argument. The tensor
\[
 T=A_0^2+A_0\wedge B_0+B_0^2
\]
satisfies $\ddc w\wedge T=B_0^3-A_0^3=0$, but $T=0$ on $M_\pm$, where $A_0=B_0=\omega_0$ has rank one. The resulting equation is not globally elliptic, even though both potentials and their common density are smooth.
